% Options for packages loaded elsewhere
% Options for packages loaded elsewhere
\PassOptionsToPackage{unicode}{hyperref}
\PassOptionsToPackage{hyphens}{url}
\PassOptionsToPackage{dvipsnames,svgnames,x11names}{xcolor}
%
\documentclass[
  letterpaper,
  DIV=11,
  numbers=noendperiod]{scrartcl}
\usepackage{xcolor}
\usepackage[margin=1in]{geometry}
\usepackage{amsmath,amssymb}
\setcounter{secnumdepth}{5}
\usepackage{iftex}
\ifPDFTeX
  \usepackage[T1]{fontenc}
  \usepackage[utf8]{inputenc}
  \usepackage{textcomp} % provide euro and other symbols
\else % if luatex or xetex
  \usepackage{unicode-math} % this also loads fontspec
  \defaultfontfeatures{Scale=MatchLowercase}
  \defaultfontfeatures[\rmfamily]{Ligatures=TeX,Scale=1}
\fi
\usepackage{lmodern}
\ifPDFTeX\else
  % xetex/luatex font selection
  \setmainfont[]{Latin Modern Roman}
  \setsansfont[]{Latin Modern Sans}
  \setmonofont[]{Latin Modern Mono}
\fi
% Use upquote if available, for straight quotes in verbatim environments
\IfFileExists{upquote.sty}{\usepackage{upquote}}{}
\IfFileExists{microtype.sty}{% use microtype if available
  \usepackage[]{microtype}
  \UseMicrotypeSet[protrusion]{basicmath} % disable protrusion for tt fonts
}{}
\makeatletter
\@ifundefined{KOMAClassName}{% if non-KOMA class
  \IfFileExists{parskip.sty}{%
    \usepackage{parskip}
  }{% else
    \setlength{\parindent}{0pt}
    \setlength{\parskip}{6pt plus 2pt minus 1pt}}
}{% if KOMA class
  \KOMAoptions{parskip=half}}
\makeatother
% Make \paragraph and \subparagraph free-standing
\makeatletter
\ifx\paragraph\undefined\else
  \let\oldparagraph\paragraph
  \renewcommand{\paragraph}{
    \@ifstar
      \xxxParagraphStar
      \xxxParagraphNoStar
  }
  \newcommand{\xxxParagraphStar}[1]{\oldparagraph*{#1}\mbox{}}
  \newcommand{\xxxParagraphNoStar}[1]{\oldparagraph{#1}\mbox{}}
\fi
\ifx\subparagraph\undefined\else
  \let\oldsubparagraph\subparagraph
  \renewcommand{\subparagraph}{
    \@ifstar
      \xxxSubParagraphStar
      \xxxSubParagraphNoStar
  }
  \newcommand{\xxxSubParagraphStar}[1]{\oldsubparagraph*{#1}\mbox{}}
  \newcommand{\xxxSubParagraphNoStar}[1]{\oldsubparagraph{#1}\mbox{}}
\fi
\makeatother

\usepackage{color}
\usepackage{fancyvrb}
\newcommand{\VerbBar}{|}
\newcommand{\VERB}{\Verb[commandchars=\\\{\}]}
\DefineVerbatimEnvironment{Highlighting}{Verbatim}{commandchars=\\\{\}}
% Add ',fontsize=\small' for more characters per line
\usepackage{framed}
\definecolor{shadecolor}{RGB}{241,243,245}
\newenvironment{Shaded}{\begin{snugshade}}{\end{snugshade}}
\newcommand{\AlertTok}[1]{\textcolor[rgb]{0.68,0.00,0.00}{#1}}
\newcommand{\AnnotationTok}[1]{\textcolor[rgb]{0.37,0.37,0.37}{#1}}
\newcommand{\AttributeTok}[1]{\textcolor[rgb]{0.40,0.45,0.13}{#1}}
\newcommand{\BaseNTok}[1]{\textcolor[rgb]{0.68,0.00,0.00}{#1}}
\newcommand{\BuiltInTok}[1]{\textcolor[rgb]{0.00,0.23,0.31}{#1}}
\newcommand{\CharTok}[1]{\textcolor[rgb]{0.13,0.47,0.30}{#1}}
\newcommand{\CommentTok}[1]{\textcolor[rgb]{0.37,0.37,0.37}{#1}}
\newcommand{\CommentVarTok}[1]{\textcolor[rgb]{0.37,0.37,0.37}{\textit{#1}}}
\newcommand{\ConstantTok}[1]{\textcolor[rgb]{0.56,0.35,0.01}{#1}}
\newcommand{\ControlFlowTok}[1]{\textcolor[rgb]{0.00,0.23,0.31}{\textbf{#1}}}
\newcommand{\DataTypeTok}[1]{\textcolor[rgb]{0.68,0.00,0.00}{#1}}
\newcommand{\DecValTok}[1]{\textcolor[rgb]{0.68,0.00,0.00}{#1}}
\newcommand{\DocumentationTok}[1]{\textcolor[rgb]{0.37,0.37,0.37}{\textit{#1}}}
\newcommand{\ErrorTok}[1]{\textcolor[rgb]{0.68,0.00,0.00}{#1}}
\newcommand{\ExtensionTok}[1]{\textcolor[rgb]{0.00,0.23,0.31}{#1}}
\newcommand{\FloatTok}[1]{\textcolor[rgb]{0.68,0.00,0.00}{#1}}
\newcommand{\FunctionTok}[1]{\textcolor[rgb]{0.28,0.35,0.67}{#1}}
\newcommand{\ImportTok}[1]{\textcolor[rgb]{0.00,0.46,0.62}{#1}}
\newcommand{\InformationTok}[1]{\textcolor[rgb]{0.37,0.37,0.37}{#1}}
\newcommand{\KeywordTok}[1]{\textcolor[rgb]{0.00,0.23,0.31}{\textbf{#1}}}
\newcommand{\NormalTok}[1]{\textcolor[rgb]{0.00,0.23,0.31}{#1}}
\newcommand{\OperatorTok}[1]{\textcolor[rgb]{0.37,0.37,0.37}{#1}}
\newcommand{\OtherTok}[1]{\textcolor[rgb]{0.00,0.23,0.31}{#1}}
\newcommand{\PreprocessorTok}[1]{\textcolor[rgb]{0.68,0.00,0.00}{#1}}
\newcommand{\RegionMarkerTok}[1]{\textcolor[rgb]{0.00,0.23,0.31}{#1}}
\newcommand{\SpecialCharTok}[1]{\textcolor[rgb]{0.37,0.37,0.37}{#1}}
\newcommand{\SpecialStringTok}[1]{\textcolor[rgb]{0.13,0.47,0.30}{#1}}
\newcommand{\StringTok}[1]{\textcolor[rgb]{0.13,0.47,0.30}{#1}}
\newcommand{\VariableTok}[1]{\textcolor[rgb]{0.07,0.07,0.07}{#1}}
\newcommand{\VerbatimStringTok}[1]{\textcolor[rgb]{0.13,0.47,0.30}{#1}}
\newcommand{\WarningTok}[1]{\textcolor[rgb]{0.37,0.37,0.37}{\textit{#1}}}

\usepackage{longtable,booktabs,array}
\usepackage{calc} % for calculating minipage widths
% Correct order of tables after \paragraph or \subparagraph
\usepackage{etoolbox}
\makeatletter
\patchcmd\longtable{\par}{\if@noskipsec\mbox{}\fi\par}{}{}
\makeatother
% Allow footnotes in longtable head/foot
\IfFileExists{footnotehyper.sty}{\usepackage{footnotehyper}}{\usepackage{footnote}}
\makesavenoteenv{longtable}
\usepackage{graphicx}
\makeatletter
\newsavebox\pandoc@box
\newcommand*\pandocbounded[1]{% scales image to fit in text height/width
  \sbox\pandoc@box{#1}%
  \Gscale@div\@tempa{\textheight}{\dimexpr\ht\pandoc@box+\dp\pandoc@box\relax}%
  \Gscale@div\@tempb{\linewidth}{\wd\pandoc@box}%
  \ifdim\@tempb\p@<\@tempa\p@\let\@tempa\@tempb\fi% select the smaller of both
  \ifdim\@tempa\p@<\p@\scalebox{\@tempa}{\usebox\pandoc@box}%
  \else\usebox{\pandoc@box}%
  \fi%
}
% Set default figure placement to htbp
\def\fps@figure{htbp}
\makeatother


% definitions for citeproc citations
\NewDocumentCommand\citeproctext{}{}
\NewDocumentCommand\citeproc{mm}{%
  \begingroup\def\citeproctext{#2}\cite{#1}\endgroup}
\makeatletter
 % allow citations to break across lines
 \let\@cite@ofmt\@firstofone
 % avoid brackets around text for \cite:
 \def\@biblabel#1{}
 \def\@cite#1#2{{#1\if@tempswa , #2\fi}}
\makeatother
\newlength{\cslhangindent}
\setlength{\cslhangindent}{1.5em}
\newlength{\csllabelwidth}
\setlength{\csllabelwidth}{3em}
\newenvironment{CSLReferences}[2] % #1 hanging-indent, #2 entry-spacing
 {\begin{list}{}{%
  \setlength{\itemindent}{0pt}
  \setlength{\leftmargin}{0pt}
  \setlength{\parsep}{0pt}
  % turn on hanging indent if param 1 is 1
  \ifodd #1
   \setlength{\leftmargin}{\cslhangindent}
   \setlength{\itemindent}{-1\cslhangindent}
  \fi
  % set entry spacing
  \setlength{\itemsep}{#2\baselineskip}}}
 {\end{list}}
\usepackage{calc}
\newcommand{\CSLBlock}[1]{\hfill\break\parbox[t]{\linewidth}{\strut\ignorespaces#1\strut}}
\newcommand{\CSLLeftMargin}[1]{\parbox[t]{\csllabelwidth}{\strut#1\strut}}
\newcommand{\CSLRightInline}[1]{\parbox[t]{\linewidth - \csllabelwidth}{\strut#1\strut}}
\newcommand{\CSLIndent}[1]{\hspace{\cslhangindent}#1}



\setlength{\emergencystretch}{3em} % prevent overfull lines

\providecommand{\tightlist}{%
  \setlength{\itemsep}{0pt}\setlength{\parskip}{0pt}}



 


% Colors and section/title styling using KOMA-Script interfaces
\usepackage{xcolor}
\definecolor{sectionblue}{HTML}{2563eb}

% KOMA: headings and title/subtitle colors
\setkomafont{title}{\color{sectionblue}\bfseries\Huge}
\setkomafont{subtitle}{\color{sectionblue}\large}
\setkomafont{section}{\color{sectionblue}\bfseries\Large}
\setkomafont{subsection}{\color{sectionblue}\bfseries\large}

% Code block styling via Shaded redefinition
\usepackage{tcolorbox}
\tcbuselibrary{skins,breakable}
\definecolor{codebg}{HTML}{F0F8FF}
\renewenvironment{Shaded}{%
  \begin{tcolorbox}[%
    enhanced,%
    colback=codebg,%
    colframe=codebg,%
    borderline west={3pt}{0pt}{sectionblue},%
    fontupper=\small\ttfamily,% reduce font size and force monospace for code
    boxrule=0pt,%
    arc=0pt,%
    boxsep=5pt,%
    left=2mm,%
    right=2mm,%
    top=2mm,%
    bottom=2mm% 
  ]% 
}{%
  \end{tcolorbox}%
}
\KOMAoption{captions}{tableheading}
\makeatletter
\@ifpackageloaded{caption}{}{\usepackage{caption}}
\AtBeginDocument{%
\ifdefined\contentsname
  \renewcommand*\contentsname{Table of contents}
\else
  \newcommand\contentsname{Table of contents}
\fi
\ifdefined\listfigurename
  \renewcommand*\listfigurename{List of Figures}
\else
  \newcommand\listfigurename{List of Figures}
\fi
\ifdefined\listtablename
  \renewcommand*\listtablename{List of Tables}
\else
  \newcommand\listtablename{List of Tables}
\fi
\ifdefined\figurename
  \renewcommand*\figurename{Figure}
\else
  \newcommand\figurename{Figure}
\fi
\ifdefined\tablename
  \renewcommand*\tablename{Table}
\else
  \newcommand\tablename{Table}
\fi
}
\@ifpackageloaded{float}{}{\usepackage{float}}
\floatstyle{ruled}
\@ifundefined{c@chapter}{\newfloat{codelisting}{h}{lop}}{\newfloat{codelisting}{h}{lop}[chapter]}
\floatname{codelisting}{Listing}
\newcommand*\listoflistings{\listof{codelisting}{List of Listings}}
\makeatother
\makeatletter
\makeatother
\makeatletter
\@ifpackageloaded{caption}{}{\usepackage{caption}}
\@ifpackageloaded{subcaption}{}{\usepackage{subcaption}}
\makeatother
\usepackage{bookmark}
\IfFileExists{xurl.sty}{\usepackage{xurl}}{} % add URL line breaks if available
\urlstyle{same}
\hypersetup{
  pdftitle={Hotel Cancellation Risk (HCR)},
  colorlinks=true,
  linkcolor={blue},
  filecolor={Maroon},
  citecolor={Blue},
  urlcolor={Blue},
  pdfcreator={LaTeX via pandoc}}


\title{Hotel Cancellation Risk (HCR)}
\usepackage{etoolbox}
\makeatletter
\providecommand{\subtitle}[1]{% add subtitle to \maketitle
  \apptocmd{\@title}{\par {\large #1 \par}}{}{}
}
\makeatother
\subtitle{Leak-Aware Prediction at Booking Time}
\author{}
\date{}
\begin{document}
\maketitle

\renewcommand*\contentsname{Table of contents}
{
\hypersetup{linkcolor=}
\setcounter{tocdepth}{3}
\tableofcontents
}

\section{Project Introduction and Problem
Description}\label{project-introduction-and-problem-description}

The goal of this project is to estimate---at the moment a hotel booking
is created---the probability that the reservation will later be
canceled.

Learning approach: supervised learning --- specifically a binary
classification task (predicting \texttt{is\_canceled} in \{0,1\}). We
treat this as a discriminative, supervised classification problem where
models learn P(cancel \textbar{} booking-time features) from labeled
historical bookings.

Algorithms (explicit categorization): - Regularized logistic regression
--- a linear, parametric, probabilistic discriminative classifier. Fast,
interpretable, and a good baseline; regularization (L1/L2) helps control
overfitting and supports feature selection/interpretability. -
Gradient-boosted trees (e.g., XGBoost/LightGBM/CatBoost) --- a
tree-based ensemble (non-parametric) built by boosting decision trees.
Captures nonlinearities and feature interactions, handles heterogeneous
feature types and missing values well, and often yields strong
predictive performance.

Rationale: using both a simple linear probabilistic model (logistic
regression) and a flexible ensemble (gradient-boosted trees) gives a
balance between interpretability and predictive power. Both are
supervised approaches; we will compare their ranking and calibration
performance and apply probability calibration where needed so outputs
are actionable for operations (overbooking rules, targeted outreach,
deposit policies).

Evaluation mirrors real use with a chronological split (earlier bookings
for training, later for validation, the most recent for testing). We
will report PR-AUC (primary) and ROC-AUC (secondary) to assess ranking
quality, and the Brier score with a reliability curve to assess
probability calibration.

Problem setup (concise): - Decision time: booking creation moment. -
Target: is\_canceled in \{0,1\}; predict P(cancel \textbar{}
booking-time features). - Evaluation sketch: time-aware split; metrics
as above; post-training calibration.

\section{Data Loading \& Initial
Inspection}\label{data-loading-initial-inspection}

For this project we specifically need an individual-level hotel booking
dataset that includes the final cancellation outcome and the fields
known at booking time. The Hotel Booking Demand dataset satisfies these
requirements.

\subsection{Dataset}\label{dataset}

We use the Hotel Booking Demand dataset, which contains reservation
records for two Portuguese hotels---one city hotel (H2) and one resort
hotel (H1)---covering arrivals from 1 July 2015 to 31 August 2017. In
total there are 119,390 bookings described by 32 variables (e.g., lead
time, market segment, distribution channel, deposit type, guest history,
stay pattern); all personally identifying information was removed by the
authors.

The dataset originates from Antonio, Almeida, and Nunes (Antonio,
Almeida, and Nunes (2019)), who exported and anonymized the data from
the hotels' operational systems; it includes both realized and canceled
bookings and was released specifically to support academic and applied
research on hotel demand.

We will use the widely mirrored Kaggle copy curated by Jesse Mostipak
(Mostipak (2020)), which repackages the original data as a single CSV
for convenient loading.

This dataset is public, well-documented, and directly aligned with our
goal of predicting cancellation risk at booking time: it contains both
outcomes (canceled vs not canceled) and rich booking-time fields, spans
multiple years and two hotel types, and is easy to reproduce from a
single CSV.

Next we will download and unzip the dataset using Kaggle API.

\subsection{Initial data inspection}\label{initial-data-inspection}

Load the dataset and show a quick overview before we begin
preprocessing: a preview of rows, the column types, and simple
descriptive statistics. If the CSV is missing the code will raise a
clear error so you can run the ingestion step first.

\begin{Shaded}
\begin{Highlighting}[]
\ImportTok{import}\NormalTok{ pandas }\ImportTok{as}\NormalTok{ pd}
\ImportTok{from}\NormalTok{ pathlib }\ImportTok{import}\NormalTok{ Path}
\ImportTok{from}\NormalTok{ IPython.display }\ImportTok{import}\NormalTok{ display, Markdown}
\NormalTok{data\_path }\OperatorTok{=} \StringTok{"../data/raw/hotel\_bookings.csv"}
\NormalTok{df }\OperatorTok{=}\NormalTok{ pd.read\_csv(data\_path)}
\end{Highlighting}
\end{Shaded}

df.info() gives the following information:

\textbf{Rows:} 119,390\\
\textbf{Columns:} 32

\begin{longtable}[]{@{}
  >{\raggedleft\arraybackslash}p{(\linewidth - 10\tabcolsep) * \real{0.0519}}
  >{\raggedright\arraybackslash}p{(\linewidth - 10\tabcolsep) * \real{0.4286}}
  >{\raggedright\arraybackslash}p{(\linewidth - 10\tabcolsep) * \real{0.1299}}
  >{\raggedleft\arraybackslash}p{(\linewidth - 10\tabcolsep) * \real{0.1299}}
  >{\raggedleft\arraybackslash}p{(\linewidth - 10\tabcolsep) * \real{0.1169}}
  >{\raggedleft\arraybackslash}p{(\linewidth - 10\tabcolsep) * \real{0.1429}}@{}}
\toprule\noalign{}
\begin{minipage}[b]{\linewidth}\raggedleft
\#
\end{minipage} & \begin{minipage}[b]{\linewidth}\raggedright
Column
\end{minipage} & \begin{minipage}[b]{\linewidth}\raggedright
dtype
\end{minipage} & \begin{minipage}[b]{\linewidth}\raggedleft
Non-null
\end{minipage} & \begin{minipage}[b]{\linewidth}\raggedleft
Missing
\end{minipage} & \begin{minipage}[b]{\linewidth}\raggedleft
Missing \%
\end{minipage} \\
\midrule\noalign{}
\endhead
\bottomrule\noalign{}
\endlastfoot
0 & hotel & object & 119,390 & 0 & 0.00\% \\
1 & is\_canceled & int64 & 119,390 & 0 & 0.00\% \\
2 & lead\_time & int64 & 119,390 & 0 & 0.00\% \\
3 & arrival\_date\_year & int64 & 119,390 & 0 & 0.00\% \\
4 & arrival\_date\_month & object & 119,390 & 0 & 0.00\% \\
5 & arrival\_date\_week\_number & int64 & 119,390 & 0 & 0.00\% \\
6 & arrival\_date\_day\_of\_month & int64 & 119,390 & 0 & 0.00\% \\
7 & stays\_in\_weekend\_nights & int64 & 119,390 & 0 & 0.00\% \\
8 & stays\_in\_week\_nights & int64 & 119,390 & 0 & 0.00\% \\
9 & adults & int64 & 119,390 & 0 & 0.00\% \\
10 & children & float64 & 119,386 & 4 & 0.00\% \\
11 & babies & int64 & 119,390 & 0 & 0.00\% \\
12 & meal & object & 119,390 & 0 & 0.00\% \\
13 & country & object & 118,902 & 488 & 0.41\% \\
14 & market\_segment & object & 119,390 & 0 & 0.00\% \\
15 & distribution\_channel & object & 119,390 & 0 & 0.00\% \\
16 & is\_repeated\_guest & int64 & 119,390 & 0 & 0.00\% \\
17 & previous\_cancellations & int64 & 119,390 & 0 & 0.00\% \\
18 & previous\_bookings\_not\_canceled & int64 & 119,390 & 0 & 0.00\% \\
19 & reserved\_room\_type & object & 119,390 & 0 & 0.00\% \\
20 & assigned\_room\_type & object & 119,390 & 0 & 0.00\% \\
21 & booking\_changes & int64 & 119,390 & 0 & 0.00\% \\
22 & deposit\_type & object & 119,390 & 0 & 0.00\% \\
23 & agent & float64 & 103,050 & 16,340 & 13.69\% \\
24 & company & float64 & 6,797 & 112,593 & 94.31\% \\
25 & days\_in\_waiting\_list & int64 & 119,390 & 0 & 0.00\% \\
26 & customer\_type & object & 119,390 & 0 & 0.00\% \\
27 & adr & float64 & 119,390 & 0 & 0.00\% \\
28 & required\_car\_parking\_spaces & int64 & 119,390 & 0 & 0.00\% \\
29 & total\_of\_special\_requests & int64 & 119,390 & 0 & 0.00\% \\
30 & reservation\_status & object & 119,390 & 0 & 0.00\% \\
31 & reservation\_status\_date & object & 119,390 & 0 & 0.00\% \\
\end{longtable}

Dataset summary (tabulated):

\begin{itemize}
\tightlist
\item
  Number of samples (rows): \textbf{119,390}
\item
  Number of features (columns): \textbf{32}
\item
  CSV file size on disk: \textbf{\textasciitilde16.9 MB}
\item
  Data provenance: single-table CSV (Kaggle mirror of the Hotel Booking
  Demand dataset).
\end{itemize}

Data types (summary): - int64: 16 - float64: 4 - object (string /
categorical): 12

Missing-value highlights (important columns):

\begin{itemize}
\tightlist
\item
  \texttt{agent}: 16,340 missing (\(\approx 13.69\%\)) --- mostly NA
  where bookings have no agent ID
\item
  \texttt{company}: 112,593 missing (\(\approx 94.31\%\)) --- most
  bookings are not associated with a company
\item
  \texttt{country}: 488 missing (\(\approx 0.41\%\))
\end{itemize}

Key features (short descriptions):

\begin{itemize}
\tightlist
\item
  \texttt{is\_canceled} (target): 0/1 indicator whether the booking was
  canceled --- the supervised label for binary classification.
\item
  \texttt{lead\_time}: number of days between booking and arrival ---
  often predictive of cancellation behaviour.
\item
  \texttt{previous\_cancellations}: historical number of cancellations
  by the same guest --- a strong signal of cancellation propensity.
\item
  \texttt{previous\_bookings\_not\_canceled}: prior non-canceled
  bookings for the guest.
\item
  \texttt{is\_repeated\_guest}: binary indicator if the guest is a
  repeat customer.
\item
  \texttt{adr}: average daily rate --- price information associated with
  the booking.
\item
  \texttt{hotel}: hotel type (city vs resort) --- two hotel sites in
  dataset.
\item
  \texttt{reservation\_status} / \texttt{reservation\_status\_date}:
  final reservation outcome and date (stored as text; parse to dates for
  temporal checks).
\item
  \texttt{booking\_changes}: number of changes made to the booking
  (edits) --- may signal instability.
\item
  \texttt{agent}, \texttt{company}: identifiers for booking agent or
  company (sparse/missing); handle carefully or create a `missing'
  category.
\item
  \texttt{arrival\_date\_*} (year/month/day): arrival date broken into
  components --- used for chronological splitting and seasonality
  features.
\item
  \texttt{total\_of\_special\_requests}, \texttt{customer\_type},
  \texttt{deposit\_type}, \texttt{market\_segment},
  \texttt{distribution\_channel}: operational/categorical fields useful
  for segmentation and modelling.
\end{itemize}

Notes on shape and sources: This dataset is a single, flattened table
(one CSV) assembled from operational exports by the dataset authors and
mirrored on Kaggle; it is not a multi-table relational dataset.

Note: children, agent, and company are float64 because of missing
values; reservation\_status\_date is stored as text and can be parsed to
a proper date type.

\section{Preprocess and Data
Cleaning}\label{preprocess-and-data-cleaning}

\subsection{Missing Value Treatment}\label{missing-value-treatment}

\textbf{Step}: Handle missing values in identified columns

\begin{itemize}
\tightlist
\item
  \texttt{children}: 4 missing (0.00\%) → Fill with 0
\item
  \texttt{country}: 488 missing (0.41\%) → needs a closer analysis
\item
  \texttt{agent}: 16,340 missing (13.69\%) → Fill with 0
\item
  \texttt{company}: 112,593 missing (94.31\%) → Fill with 0
\end{itemize}

\textbf{Why}: Missing values prevent model training.

\begin{itemize}
\tightlist
\item
  For \texttt{children}, 0 is logical default.
\item
  For \texttt{agent} and \texttt{company}, missingness indicates direct
  bookings (no intermediary), so 0 represents this business reality.
\item
  High missingness in \texttt{company} (94\%) suggests most bookings
  aren't corporate.
\item
  Country needs a closer look, which we will do in the next subsection:
\end{itemize}

\begin{Shaded}
\begin{Highlighting}[]
\NormalTok{df\_cleaned }\OperatorTok{=}\NormalTok{ df.copy()}
\NormalTok{df\_cleaned[}\StringTok{\textquotesingle{}children\textquotesingle{}}\NormalTok{] }\OperatorTok{=}\NormalTok{ df\_cleaned[}\StringTok{\textquotesingle{}children\textquotesingle{}}\NormalTok{].fillna(}\DecValTok{0}\NormalTok{)}
\NormalTok{df\_cleaned[}\StringTok{\textquotesingle{}agent\textquotesingle{}}\NormalTok{] }\OperatorTok{=}\NormalTok{ df\_cleaned[}\StringTok{\textquotesingle{}agent\textquotesingle{}}\NormalTok{].fillna(}\DecValTok{0}\NormalTok{)}
\NormalTok{df\_cleaned[}\StringTok{\textquotesingle{}company\textquotesingle{}}\NormalTok{] }\OperatorTok{=}\NormalTok{ df\_cleaned[}\StringTok{\textquotesingle{}company\textquotesingle{}}\NormalTok{].fillna(}\DecValTok{0}\NormalTok{)}
\end{Highlighting}
\end{Shaded}

\subsubsection{Missing Values - Country
Analysis}\label{missing-values---country-analysis}

After doing some investigation, missing country bookings show
dramatically different patterns:

\begin{itemize}
\tightlist
\item
  13.7\% cancellation rate vs 37.1\% overall
\item
  95\% resort hotels vs 34\% overall\\
\item
  38\% corporate distribution vs 6\% overall
\item
  32-day lead time vs 104-day overall
\item
  Likely represents corporate/business segment with distinct booking
  behavior
\end{itemize}

\textbf{Why}: Missing country correlates with a low-risk, corporate
booking segment. The `Unknown' category captures valuable predictive
signal about booking certainty and should be preserved for modeling.

\textbf{Step}: Retain missing country values as `Unknown' category

\begin{itemize}
\tightlist
\item
  Fill 488 missing \texttt{country} values with `Unknown'
\item
  Treat as legitimate business segment for modeling
\end{itemize}

\begin{Shaded}
\begin{Highlighting}[]
\NormalTok{df\_cleaned[}\StringTok{\textquotesingle{}country\textquotesingle{}}\NormalTok{] }\OperatorTok{=}\NormalTok{ df\_cleaned[}\StringTok{\textquotesingle{}country\textquotesingle{}}\NormalTok{].fillna(}\StringTok{\textquotesingle{}Unknown\textquotesingle{}}\NormalTok{)}
\end{Highlighting}
\end{Shaded}

\subsection{Data Leakage Removal}\label{data-leakage-removal}

\textbf{Step}: Remove features that reveal the target or aren't
available at booking time

\begin{itemize}
\tightlist
\item
  Remove \texttt{reservation\_status} (directly indicates cancellation)
\item
  Remove \texttt{reservation\_status\_date} (cancellation/checkout date)
\item
  Remove \texttt{deposit\_type} (During EDA, we discovered that
  deposit\_type contains post-outcome information. The ``Non Refund''
  category shows 99.4\% cancellation rate (14,494 of 14,587 bookings), )
\end{itemize}

\textbf{Why}: These features create data leakage - they contain
information that would only be known after the prediction needs to be
made. Using them would artificially inflate model performance and make
it useless in production.

We also remove 7 bookings with `Undefined' market segments or
distribution channels - these represent early data collection errors
from 2015 rather than meaningful unknown categories. This removal
affects less than 0.006\% of the dataset while eliminating anomalous
patterns (also discovered during EDA)

\begin{Shaded}
\begin{Highlighting}[]
\NormalTok{df\_cleaned }\OperatorTok{=}\NormalTok{ df\_cleaned.drop(columns}\OperatorTok{=}\NormalTok{[}\StringTok{\textquotesingle{}reservation\_status\textquotesingle{}}\NormalTok{, }\StringTok{\textquotesingle{}reservation\_status\_date\textquotesingle{}}\NormalTok{, }\StringTok{\textquotesingle{}deposit\_type\textquotesingle{}}\NormalTok{])}

\CommentTok{\# Remove undefined entries}
\NormalTok{df\_cleaned }\OperatorTok{=}\NormalTok{ df\_cleaned[}
    \OperatorTok{\textasciitilde{}}\NormalTok{df\_cleaned[}\StringTok{\textquotesingle{}market\_segment\textquotesingle{}}\NormalTok{].isin([}\StringTok{\textquotesingle{}Undefined\textquotesingle{}}\NormalTok{]) }\OperatorTok{\&} 
    \OperatorTok{\textasciitilde{}}\NormalTok{df\_cleaned[}\StringTok{\textquotesingle{}distribution\_channel\textquotesingle{}}\NormalTok{].isin([}\StringTok{\textquotesingle{}Undefined\textquotesingle{}}\NormalTok{])}
\NormalTok{]}
\end{Highlighting}
\end{Shaded}

\subsection{Outlier Detection and
Treatment}\label{outlier-detection-and-treatment}

\textbf{Step}: Clean ADR (Average Daily Rate) outliers using
business-informed approach

Average Daily Rate shows extreme values requiring strategic treatment
rather than blanket removal:

\begin{Shaded}
\begin{Highlighting}[]
\ImportTok{import}\NormalTok{ numpy }\ImportTok{as}\NormalTok{ np}

\NormalTok{adr }\OperatorTok{=}\NormalTok{ df[}\StringTok{\textquotesingle{}adr\textquotesingle{}}\NormalTok{]}
\NormalTok{q1, q3 }\OperatorTok{=}\NormalTok{ adr.quantile([}\FloatTok{0.25}\NormalTok{, }\FloatTok{0.75}\NormalTok{])}

\BuiltInTok{print}\NormalTok{(}\SpecialStringTok{f"ADR Stats (n=}\SpecialCharTok{\{}\BuiltInTok{len}\NormalTok{(adr)}\SpecialCharTok{:,\}}\SpecialStringTok{)"}\NormalTok{)}
\BuiltInTok{print}\NormalTok{(}\SpecialStringTok{f"Range: $}\SpecialCharTok{\{}\NormalTok{adr}\SpecialCharTok{.}\BuiltInTok{min}\NormalTok{()}\SpecialCharTok{:.2f\}}\SpecialStringTok{ {-} $}\SpecialCharTok{\{}\NormalTok{adr}\SpecialCharTok{.}\BuiltInTok{max}\NormalTok{()}\SpecialCharTok{:,.0f\}}\SpecialStringTok{"}\NormalTok{)}
\BuiltInTok{print}\NormalTok{(}\SpecialStringTok{f"Mean: $}\SpecialCharTok{\{}\NormalTok{adr}\SpecialCharTok{.}\NormalTok{mean()}\SpecialCharTok{:.2f\}}\SpecialStringTok{ | Median: $}\SpecialCharTok{\{}\NormalTok{adr}\SpecialCharTok{.}\NormalTok{median()}\SpecialCharTok{:.2f\}}\SpecialStringTok{"}\NormalTok{)}
\BuiltInTok{print}\NormalTok{(}\SpecialStringTok{f"1\%: $}\SpecialCharTok{\{}\NormalTok{adr}\SpecialCharTok{.}\NormalTok{quantile(}\FloatTok{0.01}\NormalTok{)}\SpecialCharTok{:.2f\}}\SpecialStringTok{ | 99,99\%: $}\SpecialCharTok{\{}\NormalTok{adr}\SpecialCharTok{.}\NormalTok{quantile(}\FloatTok{0.9999}\NormalTok{)}\SpecialCharTok{:.2f\}}\SpecialStringTok{"}\NormalTok{)}
\BuiltInTok{print}\NormalTok{(}\SpecialStringTok{f"Issues: Neg:}\SpecialCharTok{\{}\NormalTok{(adr}\OperatorTok{\textless{}}\DecValTok{0}\NormalTok{)}\SpecialCharTok{.}\BuiltInTok{sum}\NormalTok{()}\SpecialCharTok{\}}\SpecialStringTok{ Zero:}\SpecialCharTok{\{}\NormalTok{(adr}\OperatorTok{==}\DecValTok{0}\NormalTok{)}\SpecialCharTok{.}\BuiltInTok{sum}\NormalTok{()}\SpecialCharTok{\}}\SpecialStringTok{ \textgreater{}$500:}\SpecialCharTok{\{}\NormalTok{(adr}\OperatorTok{\textgreater{}}\DecValTok{500}\NormalTok{)}\SpecialCharTok{.}\BuiltInTok{sum}\NormalTok{()}\SpecialCharTok{\}}\SpecialStringTok{"}\NormalTok{)}
\NormalTok{outliers }\OperatorTok{=}\NormalTok{ (adr }\OperatorTok{\textgreater{}}\NormalTok{ q3 }\OperatorTok{+} \FloatTok{1.5}\OperatorTok{*}\NormalTok{(q3}\OperatorTok{{-}}\NormalTok{q1)).}\BuiltInTok{sum}\NormalTok{()}
\BuiltInTok{print}\NormalTok{(}\SpecialStringTok{f"IQR outliers: }\SpecialCharTok{\{}\NormalTok{outliers}\SpecialCharTok{:,\}}\SpecialStringTok{"}\NormalTok{)}
\ControlFlowTok{if}\NormalTok{ (adr }\OperatorTok{\textgreater{}} \DecValTok{1000}\NormalTok{).}\BuiltInTok{sum}\NormalTok{() }\OperatorTok{\textless{}=} \DecValTok{5}\NormalTok{: }
    \BuiltInTok{print}\NormalTok{(}\SpecialStringTok{f"\textgreater{}$1K: }\SpecialCharTok{\{}\BuiltInTok{list}\NormalTok{(adr[adr}\OperatorTok{\textgreater{}}\DecValTok{1000}\NormalTok{])}\SpecialCharTok{\}}\SpecialStringTok{"}\NormalTok{)}
\end{Highlighting}
\end{Shaded}

\begin{verbatim}
ADR Stats (n=119,390)
Range: $-6.38 - $5,400
Mean: $101.83 | Median: $94.58
1%: $0.00 | 99,99%: $388.00
Issues: Neg:1 Zero:1959 >$500:3
IQR outliers: 3,793
>$1K: [5400.0]
\end{verbatim}

\begin{itemize}
\tightlist
\item
  \textbf{Remove 1 negative ADR} (-\$6.38) - impossible hotel rate
\item
  \textbf{Remove 715 incomplete bookings} (zero ADR + zero nights) -
  data entry errors\\
\item
  \textbf{Preserve 1,244 complementary stays} (zero ADR + actual nights)
  - legitimate complimentary bookings
\item
  \textbf{Cap extreme outlier} (\(\approx 500\,€\) at 99,999th
  percentile) - likely data entry error
\end{itemize}

\textbf{Why}: Investigation revealed complementary stays represent a
distinct low-risk segment with 14.3\% cancellation rate versus 37.1\%
overall. Rather than removing all zero-price bookings, we preserve this
valuable business segment while eliminating only clearly invalid
records. This approach maintains predictive signal while ensuring data
quality.

\begin{Shaded}
\begin{Highlighting}[]
\NormalTok{df\_cleaned }\OperatorTok{=}\NormalTok{ df\_cleaned[df\_cleaned[}\StringTok{\textquotesingle{}adr\textquotesingle{}}\NormalTok{] }\OperatorTok{\textgreater{}=} \DecValTok{0}\NormalTok{]  }\CommentTok{\# Remove negative ADR}
\NormalTok{cap\_value }\OperatorTok{=}\NormalTok{ df[}\StringTok{\textquotesingle{}adr\textquotesingle{}}\NormalTok{].quantile(}\FloatTok{0.9999}\NormalTok{)}
\NormalTok{df\_cleaned[}\StringTok{\textquotesingle{}adr\textquotesingle{}}\NormalTok{] }\OperatorTok{=}\NormalTok{ df\_cleaned[}\StringTok{\textquotesingle{}adr\textquotesingle{}}\NormalTok{].clip(upper}\OperatorTok{=}\NormalTok{cap\_value)}
\end{Highlighting}
\end{Shaded}

\begin{Shaded}
\begin{Highlighting}[]
\ImportTok{import}\NormalTok{ matplotlib.pyplot }\ImportTok{as}\NormalTok{ plt}
\CommentTok{\# Histogram with linear x{-}axis and log y{-}axis}
\NormalTok{fig, ax }\OperatorTok{=}\NormalTok{ plt.subplots(figsize}\OperatorTok{=}\NormalTok{(}\DecValTok{12}\NormalTok{, }\DecValTok{6}\NormalTok{))}
\NormalTok{ax.hist(df\_cleaned[}\StringTok{\textquotesingle{}adr\textquotesingle{}}\NormalTok{], bins}\OperatorTok{=}\DecValTok{200}\NormalTok{, edgecolor}\OperatorTok{=}\StringTok{\textquotesingle{}black\textquotesingle{}}\NormalTok{, alpha}\OperatorTok{=}\FloatTok{0.7}\NormalTok{, color}\OperatorTok{=}\StringTok{\textquotesingle{}lightblue\textquotesingle{}}\NormalTok{)}
\NormalTok{ax.set\_yscale(}\StringTok{\textquotesingle{}log\textquotesingle{}}\NormalTok{)  }\CommentTok{\# Log scale on y{-}axis only}
\NormalTok{ax.set\_title(}\StringTok{\textquotesingle{}ADR Distribution {-} ALL DATA (Log Y{-}axis)\textquotesingle{}}\NormalTok{)}
\NormalTok{ax.set\_xlabel(}\StringTok{\textquotesingle{}Average Daily Rate ($)\textquotesingle{}}\NormalTok{)}
\NormalTok{ax.set\_ylabel(}\StringTok{\textquotesingle{}Frequency (Log Scale)\textquotesingle{}}\NormalTok{)}
\NormalTok{ax.grid(}\VariableTok{True}\NormalTok{, alpha}\OperatorTok{=}\FloatTok{0.3}\NormalTok{)}

\CommentTok{\# Add title with key stats}
\NormalTok{ax.set\_title(}\SpecialStringTok{f\textquotesingle{}AFTER CLEANING {-} ADR Distribution (Log Y{-}axis)}\CharTok{\textbackslash{}n}\SpecialStringTok{\textquotesingle{}}\NormalTok{)}

\NormalTok{plt.tight\_layout()}
\NormalTok{plt.show()}
\end{Highlighting}
\end{Shaded}

\pandocbounded{\includegraphics[keepaspectratio]{main_files/figure-pdf/cell-8-output-1.png}}

\begin{Shaded}
\begin{Highlighting}[]
\NormalTok{df\_cleaned }\OperatorTok{=}\NormalTok{ df\_cleaned.drop(columns}\OperatorTok{=}\NormalTok{[}\StringTok{\textquotesingle{}arrival\_date\_year\textquotesingle{}}\NormalTok{, }\StringTok{\textquotesingle{}arrival\_date\_month\textquotesingle{}}\NormalTok{, }\StringTok{\textquotesingle{}arrival\_date\_day\_of\_month\textquotesingle{}}\NormalTok{, }\StringTok{\textquotesingle{}arrival\_date\_week\_number\textquotesingle{}}\NormalTok{])}
\end{Highlighting}
\end{Shaded}

We dropped temporal features (year, month, day) since these are
arbitrary numeric values (31, 23, 11, etc.) without meaningful
predictive power for our ML model. While we could have feature
engineered seasonal variables (e.g., winter, summer) or other time-based
patterns, our dataset only spans 2 years, which is insufficient to
reliably capture or validate such temporal relationships.

\section{Exploratory Data Analysis}\label{exploratory-data-analysis}

After cleaning our data, we now conduct a thorough exploratory analysis
to understand the relationships between features and cancellation
behavior. This analysis will guide our modeling decisions and reveal
actionable business insights.

\subsection{Target Variable
Distribution}\label{target-variable-distribution}

\textbf{Purpose}: Understanding the base rate of cancellations is
crucial for setting performance benchmarks and identifying class
imbalance issues that may affect model training.

\begin{Shaded}
\begin{Highlighting}[]
\BuiltInTok{print}\NormalTok{(df\_cleaned[}\StringTok{\textquotesingle{}is\_canceled\textquotesingle{}}\NormalTok{].value\_counts(normalize}\OperatorTok{=}\VariableTok{True}\NormalTok{).}\BuiltInTok{round}\NormalTok{(}\DecValTok{3}\NormalTok{))}
\end{Highlighting}
\end{Shaded}

\begin{verbatim}
is_canceled
0    0.63
1    0.37
Name: proportion, dtype: float64
\end{verbatim}

\textbf{Finding}: With 37\% of bookings canceled, we have a moderately
imbalanced dataset. This is sufficient minority class representation for
standard algorithms, though we may benefit from stratified sampling
during model validation.

\subsection{Numerical Features
Analysis}\label{numerical-features-analysis}

\subsubsection{Correlation Analysis}\label{correlation-analysis}

\textbf{Purpose:} Correlation analysis helps identify linear
relationships between numerical features and our target variable. Strong
correlations suggest predictive power, while correlations between
features may indicate multicollinearity issues for linear models.

\begin{Shaded}
\begin{Highlighting}[]
\ImportTok{import}\NormalTok{ seaborn }\ImportTok{as}\NormalTok{ sns}
\NormalTok{numerical\_cols }\OperatorTok{=}\NormalTok{ df\_cleaned.select\_dtypes(include}\OperatorTok{=}\NormalTok{[}\StringTok{\textquotesingle{}int64\textquotesingle{}}\NormalTok{, }\StringTok{\textquotesingle{}float64\textquotesingle{}}\NormalTok{]).columns}
\NormalTok{correlation\_matrix }\OperatorTok{=}\NormalTok{ df\_cleaned[numerical\_cols].corr()}

\NormalTok{target\_corr }\OperatorTok{=}\NormalTok{ correlation\_matrix[}\StringTok{\textquotesingle{}is\_canceled\textquotesingle{}}\NormalTok{].sort\_values(ascending}\OperatorTok{=}\VariableTok{False}\NormalTok{)}
\NormalTok{plt.figure(figsize}\OperatorTok{=}\NormalTok{(}\DecValTok{8}\NormalTok{, }\DecValTok{10}\NormalTok{))}
\NormalTok{target\_corr[}\DecValTok{1}\NormalTok{:].plot(kind}\OperatorTok{=}\StringTok{\textquotesingle{}barh\textquotesingle{}}\NormalTok{)}
\NormalTok{plt.title(}\StringTok{\textquotesingle{}Feature Correlations with Cancellation\textquotesingle{}}\NormalTok{)}
\NormalTok{plt.xlabel(}\StringTok{\textquotesingle{}Correlation Coefficient\textquotesingle{}}\NormalTok{)}
\end{Highlighting}
\end{Shaded}

\begin{verbatim}
Text(0.5, 0, 'Correlation Coefficient')
\end{verbatim}

\pandocbounded{\includegraphics[keepaspectratio]{main_files/figure-pdf/cell-11-output-2.png}}

\subsubsection{Correlation Analysis --- Top
signals}\label{correlation-analysis-top-signals}

\begin{itemize}
\tightlist
\item
  \texttt{lead\_time} \textbf{(+0.29)} --- Longer lead times → higher
  cancellation risk (more time for plans to change).
\item
  \texttt{total\_special\_requests} \textbf{(−0.23)} --- Guests
  investing effort/requests are less likely to cancel.
\item
  \texttt{required\_car\_parking\_spaces} \textbf{(−0.20)} --- Planning
  to arrive by car signals commitment.
\item
  \texttt{previous\_cancellations} \textbf{(+0.11)} --- Past behavior
  predicts future behavior.
\end{itemize}

\begin{quote}
Note: Correlation is not equal causation. Effects may sharpen or shift
once we control for other variables.
\end{quote}

\subsection{Categorical Features
Analysis}\label{categorical-features-analysis}

\subsubsection{Statistical Tests}\label{statistical-tests}

\textbf{Purpose:} Quantify the predictive strength of categorical
features using Chi-square tests and Cramer's V effect sizes. This
statistical foundation guides which features to prioritize in modeling.

\begin{Shaded}
\begin{Highlighting}[]
\ImportTok{from}\NormalTok{ scipy }\ImportTok{import}\NormalTok{ stats}

\CommentTok{\# Chi{-}square test for categorical features}
\NormalTok{categorical\_cols }\OperatorTok{=}\NormalTok{ [}\StringTok{\textquotesingle{}hotel\textquotesingle{}}\NormalTok{, }\StringTok{\textquotesingle{}meal\textquotesingle{}}\NormalTok{, }\StringTok{\textquotesingle{}market\_segment\textquotesingle{}}\NormalTok{, }\StringTok{\textquotesingle{}distribution\_channel\textquotesingle{}}\NormalTok{, }
                    \StringTok{\textquotesingle{}is\_repeated\_guest\textquotesingle{}}\NormalTok{, }\StringTok{\textquotesingle{}customer\_type\textquotesingle{}}\NormalTok{, }\StringTok{\textquotesingle{}deposit\_type\textquotesingle{}}\NormalTok{]}

\NormalTok{chi2\_results }\OperatorTok{=}\NormalTok{ []}
\ControlFlowTok{for}\NormalTok{ feature }\KeywordTok{in}\NormalTok{ categorical\_cols:}
    \CommentTok{\# Create contingency table}
\NormalTok{    contingency }\OperatorTok{=}\NormalTok{ pd.crosstab(df[feature], df[}\StringTok{\textquotesingle{}is\_canceled\textquotesingle{}}\NormalTok{])}
\NormalTok{    chi2, p\_value, dof, expected }\OperatorTok{=}\NormalTok{ stats.chi2\_contingency(contingency)}
    
    \CommentTok{\# Calculate Cramér\textquotesingle{}s V for effect size}
\NormalTok{    n }\OperatorTok{=}\NormalTok{ contingency.}\BuiltInTok{sum}\NormalTok{().}\BuiltInTok{sum}\NormalTok{()}
\NormalTok{    cramers\_v }\OperatorTok{=}\NormalTok{ np.sqrt(chi2 }\OperatorTok{/}\NormalTok{ (n }\OperatorTok{*}\NormalTok{ (}\BuiltInTok{min}\NormalTok{(contingency.shape) }\OperatorTok{{-}} \DecValTok{1}\NormalTok{)))}
    
\NormalTok{    chi2\_results.append(\{}
        \StringTok{\textquotesingle{}Feature\textquotesingle{}}\NormalTok{: feature,}
        \StringTok{\textquotesingle{}Chi2\textquotesingle{}}\NormalTok{: chi2,}
        \StringTok{\textquotesingle{}p{-}value\textquotesingle{}}\NormalTok{: p\_value,}
        \StringTok{\textquotesingle{}Cramers\_V\textquotesingle{}}\NormalTok{: cramers\_v,}
        \StringTok{\textquotesingle{}Significance\textquotesingle{}}\NormalTok{: }\StringTok{\textquotesingle{}Yes\textquotesingle{}} \ControlFlowTok{if}\NormalTok{ p\_value }\OperatorTok{\textless{}} \FloatTok{0.001} \ControlFlowTok{else} \StringTok{\textquotesingle{}No\textquotesingle{}}
\NormalTok{    \})}

\NormalTok{chi2\_df }\OperatorTok{=}\NormalTok{ pd.DataFrame(chi2\_results).sort\_values(}\StringTok{\textquotesingle{}Cramers\_V\textquotesingle{}}\NormalTok{, ascending}\OperatorTok{=}\VariableTok{False}\NormalTok{)}
\BuiltInTok{print}\NormalTok{(}\StringTok{"Feature Importance (Statistical Tests)"}\NormalTok{)}
\BuiltInTok{print}\NormalTok{(chi2\_df.to\_string(index}\OperatorTok{=}\VariableTok{False}\NormalTok{))}
\end{Highlighting}
\end{Shaded}

\subsubsection{Cancellation Pattern
Analysis}\label{cancellation-pattern-analysis}

\textbf{Purpose}: Different customer segments and booking types exhibit
distinct cancellation behaviors. Understanding these patterns helps
hotels tailor policies (deposits, confirmations, overbooking) to
specific segments. We will also create lead \_time categories since the
correlation analysis showed that lead\_time has a high cancellation
risk.

\begin{Shaded}
\begin{Highlighting}[]
\CommentTok{\# Create a lead time category feature first}
\NormalTok{df\_cleaned[}\StringTok{\textquotesingle{}lead\_time\_category\textquotesingle{}}\NormalTok{] }\OperatorTok{=}\NormalTok{ pd.cut(df\_cleaned[}\StringTok{\textquotesingle{}lead\_time\textquotesingle{}}\NormalTok{], }
\NormalTok{                                           bins}\OperatorTok{=}\NormalTok{[}\DecValTok{0}\NormalTok{, }\DecValTok{7}\NormalTok{, }\DecValTok{30}\NormalTok{, }\DecValTok{90}\NormalTok{, }\DecValTok{400}\NormalTok{], }
\NormalTok{                                           labels}\OperatorTok{=}\NormalTok{[}\StringTok{\textquotesingle{}Last Minute\textquotesingle{}}\NormalTok{, }\StringTok{\textquotesingle{}Short\textquotesingle{}}\NormalTok{, }\StringTok{\textquotesingle{}Medium\textquotesingle{}}\NormalTok{, }\StringTok{\textquotesingle{}Long\textquotesingle{}}\NormalTok{])}

\NormalTok{categorical\_features }\OperatorTok{=}\NormalTok{ \{}
    \StringTok{\textquotesingle{}hotel\textquotesingle{}}\NormalTok{: }\StringTok{\textquotesingle{}Hotel Type\textquotesingle{}}\NormalTok{,}
    \StringTok{\textquotesingle{}customer\_type\textquotesingle{}}\NormalTok{: }\StringTok{\textquotesingle{}Customer Type\textquotesingle{}}\NormalTok{,}
    \StringTok{\textquotesingle{}market\_segment\textquotesingle{}}\NormalTok{: }\StringTok{\textquotesingle{}Market Segment\textquotesingle{}}\NormalTok{,}
    \StringTok{\textquotesingle{}distribution\_channel\textquotesingle{}}\NormalTok{: }\StringTok{\textquotesingle{}Distribution Channel\textquotesingle{}}\NormalTok{,}
    \StringTok{\textquotesingle{}is\_repeated\_guest\textquotesingle{}}\NormalTok{: }\StringTok{\textquotesingle{}Repeat Guest Status\textquotesingle{}}\NormalTok{,}
    \StringTok{\textquotesingle{}lead\_time\_category\textquotesingle{}}\NormalTok{: }\StringTok{\textquotesingle{}Booking Lead Time\textquotesingle{}}  \CommentTok{\# Add this}
\NormalTok{\}}

\NormalTok{fig, axes }\OperatorTok{=}\NormalTok{ plt.subplots(}\DecValTok{2}\NormalTok{, }\DecValTok{3}\NormalTok{, figsize}\OperatorTok{=}\NormalTok{(}\DecValTok{15}\NormalTok{, }\DecValTok{10}\NormalTok{))}
\NormalTok{axes }\OperatorTok{=}\NormalTok{ axes.ravel()}

\ControlFlowTok{for}\NormalTok{ idx, (feature, title) }\KeywordTok{in} \BuiltInTok{enumerate}\NormalTok{(categorical\_features.items()):}
    \CommentTok{\# Calculate cancellation rate by category}
\NormalTok{    cancel\_rate }\OperatorTok{=}\NormalTok{ df\_cleaned.groupby(feature)[}\StringTok{\textquotesingle{}is\_canceled\textquotesingle{}}\NormalTok{].agg([}\StringTok{\textquotesingle{}mean\textquotesingle{}}\NormalTok{, }\StringTok{\textquotesingle{}count\textquotesingle{}}\NormalTok{])}
\NormalTok{    cancel\_rate }\OperatorTok{=}\NormalTok{ cancel\_rate.sort\_values(}\StringTok{\textquotesingle{}mean\textquotesingle{}}\NormalTok{, ascending}\OperatorTok{=}\VariableTok{False}\NormalTok{)}
    
    \CommentTok{\# Plot}
\NormalTok{    cancel\_rate[}\StringTok{\textquotesingle{}mean\textquotesingle{}}\NormalTok{].plot(kind}\OperatorTok{=}\StringTok{\textquotesingle{}bar\textquotesingle{}}\NormalTok{, ax}\OperatorTok{=}\NormalTok{axes[idx], color}\OperatorTok{=}\StringTok{\textquotesingle{}steelblue\textquotesingle{}}\NormalTok{)}
\NormalTok{    axes[idx].set\_title(}\SpecialStringTok{f\textquotesingle{}Cancellation Rate by }\SpecialCharTok{\{}\NormalTok{title}\SpecialCharTok{\}}\SpecialStringTok{\textquotesingle{}}\NormalTok{)}
\NormalTok{    axes[idx].set\_ylabel(}\StringTok{\textquotesingle{}Cancellation Rate\textquotesingle{}}\NormalTok{)}
\NormalTok{    axes[idx].axhline(y}\OperatorTok{=}\FloatTok{0.37}\NormalTok{, color}\OperatorTok{=}\StringTok{\textquotesingle{}red\textquotesingle{}}\NormalTok{, linestyle}\OperatorTok{=}\StringTok{\textquotesingle{}{-}{-}\textquotesingle{}}\NormalTok{, alpha}\OperatorTok{=}\FloatTok{0.5}\NormalTok{, label}\OperatorTok{=}\StringTok{\textquotesingle{}Overall Rate\textquotesingle{}}\NormalTok{)}
\NormalTok{    axes[idx].legend()}
    
    \CommentTok{\# Add sample size as text}
    \ControlFlowTok{for}\NormalTok{ i, (index, row) }\KeywordTok{in} \BuiltInTok{enumerate}\NormalTok{(cancel\_rate.iterrows()):}
\NormalTok{        axes[idx].text(i, row[}\StringTok{\textquotesingle{}mean\textquotesingle{}}\NormalTok{] }\OperatorTok{+} \FloatTok{0.01}\NormalTok{, }\SpecialStringTok{f\textquotesingle{}n=}\SpecialCharTok{\{}\NormalTok{row[}\StringTok{"count"}\NormalTok{]}\SpecialCharTok{:,\}}\SpecialStringTok{\textquotesingle{}}\NormalTok{, }
\NormalTok{                      ha}\OperatorTok{=}\StringTok{\textquotesingle{}center\textquotesingle{}}\NormalTok{, va}\OperatorTok{=}\StringTok{\textquotesingle{}bottom\textquotesingle{}}\NormalTok{, fontsize}\OperatorTok{=}\DecValTok{8}\NormalTok{)}

\NormalTok{plt.suptitle(}\StringTok{\textquotesingle{}Cancellation Patterns Across Customer Segments\textquotesingle{}}\NormalTok{, fontsize}\OperatorTok{=}\DecValTok{14}\NormalTok{, y}\OperatorTok{=}\FloatTok{1.02}\NormalTok{)}
\NormalTok{plt.tight\_layout()}
\end{Highlighting}
\end{Shaded}

\pandocbounded{\includegraphics[keepaspectratio]{main_files/figure-pdf/cell-13-output-1.png}}

\subsection{Feature Importance
Summary}\label{feature-importance-summary}

Based on statistical tests and correlation analysis, we prioritize
features for modeling:

\subsubsection{Feature Ranking for
Modeling}\label{feature-ranking-for-modeling}

\textbf{Must Include (Effect Size \textgreater{} 0.15):}

\begin{itemize}
\tightlist
\item
  market\_segment (V=0.27)
\item
  distribution\_channel (V=0.18)
\item
  lead\_time (r=0.29)
\end{itemize}

\textbf{Should Include (Effect Size 0.05-0.15):}

\begin{itemize}
\tightlist
\item
  hotel (V=0.14)
\item
  customer\_type (V=0.14)
\item
  is\_repeated\_guest (V=0.08)\emph{ }Keep despite low effect - business
  critical
\end{itemize}

\textbf{Consider Dropping:}

\begin{itemize}
\tightlist
\item
  meal (V=0.05) - minimal predictive value
\end{itemize}

\textbf{Feature Engineering Opportunities:}

\begin{itemize}
\tightlist
\item
  lead\_time bins (clear threshold effects)
\item
  market\_segment × distribution interactions
\item
  total\_stay\_nights (weekend + week)
\end{itemize}

\subsection{Modelling Implications}\label{modelling-implications}

Our EDA highlights a mix of linear and non-linear relationships.
Logistic regression offers a transparent baseline, while
gradient-boosted trees (GBT) can capture the strong interactions (e.g.,
market segment × channel) and non-linear thresholds in lead time.

Key engineered features include binned lead\_time, total\_stay\_nights,
and selected interaction terms (especially market\_segment ×
distribution\_channel). Customer commitment signals such as special
requests and parking spaces should also be retained.

The target variable shows a moderate imbalance (37\% cancellations).
Standard classifiers can handle this, but stratified splits and
calibration methods (e.g., class weights, probability calibration) will
help ensure robust probability estimates.

Together, these insights motivate a two-model strategy: a simple,
interpretable logistic regression baseline and a flexible,
high-performance GBT model that can leverage the richer feature set.

\subsection{Finalized Feature
Engineering}\label{finalized-feature-engineering}

Running \texttt{src/hcr/features.py} on the cleaned bookings produces
\texttt{artifacts/features.csv} with \textbf{50 modeling features} plus
the target column. Every column is numeric: raw signals (lead\_time,
adr, guest counts), engineered aggregations (total\_stay\_nights),
lead\_time buckets, and curated segment×channel interactions with rare
pairings grouped into \texttt{Other}.

This table is our canonical modeling input. Logistic regression can
consume it after optional standardization of the continuous fields,
while tree-based models (e.g., gradient-boosted trees) can use it
directly. Regenerating the file is deterministic via
\texttt{PYTHONPATH=src\ python3\ -m\ hcr.features}, ensuring training
reproducibility.

\section{Models}\label{models}

We train three classifiers on \texttt{artifacts/features.csv} to balance
interpretability, predictive power, and baseline comparison. All operate
on the same engineered feature table described above:

\begin{itemize}
\item
  \textbf{Regularized Logistic Regression}: Baseline reference model.
  Before fitting we standardize continuous fields (lead\_time, adr,
  total\_stay\_nights, etc.) and use class-weight balancing so the 37\%
  cancellation rate does not skew the decision boundary. The model
  yields calibrated probabilities and feature weights for stakeholder
  interpretation.
\item
  \textbf{Gradient-Boosted Trees (GBT)}: Non-linear learner (e.g.,
  LightGBM/XGBoost) tuned for shallow depth, moderate learning rate, and
  early stopping. GBT naturally handles the one-hot encoded interactions
  and lead\_time thresholds, capturing non-additive effects highlighted
  in the EDA. Class weights or scale\_pos\_weight mirror the imbalance
  handling.
\item
  \textbf{Random Forest (baseline tree ensemble)}: Serves as a stable
  bagging-based baseline against which boosted methods can be compared.
  We tune n\_estimators and max\_depth to control variance and runtime,
  and use class\_weight balancing to handle the moderate cancellation
  imbalance.
\end{itemize}

All three models are evaluated under the time-aware split defined in the
training pipeline, with PR-AUC as the primary metric, ROC-AUC as
secondary, and calibration curves to validate probability quality.

\begin{longtable}[]{@{}
  >{\raggedright\arraybackslash}p{(\linewidth - 4\tabcolsep) * \real{0.1296}}
  >{\raggedright\arraybackslash}p{(\linewidth - 4\tabcolsep) * \real{0.4630}}
  >{\raggedright\arraybackslash}p{(\linewidth - 4\tabcolsep) * \real{0.4074}}@{}}
\toprule\noalign{}
\begin{minipage}[b]{\linewidth}\raggedright
Model
\end{minipage} & \begin{minipage}[b]{\linewidth}\raggedright
Parameters \& Features
\end{minipage} & \begin{minipage}[b]{\linewidth}\raggedright
Key Characteristics
\end{minipage} \\
\midrule\noalign{}
\endhead
\bottomrule\noalign{}
\endlastfoot
\textbf{Logistic Regression} & L1/L2 penalty grid: \{0.01, 0.1, 1,
10\}Standardized features + one-hot encoding & Transparent baseline with
interpretable coefficientsClass-balanced for 37\% cancellation rate \\
\textbf{XGBoost} & Learning rate \{0.05, 0.1\}, max\_depth \{3, 5,
7\}All 50 engineered features (bins, interactions) & Captures non-linear
patterns \& feature interactionsScale\_pos\_weight ≈ 1.7 for imbalance
handling \\
\textbf{Random Forest} & n\_estimators \{100, 300\}Same feature set as
XGBoost & Stable bagging-based ensemble baselineClass-balanced with
robust tree averaging \\
\end{longtable}

\section{Training}\label{training}

\subsection{Training Results}\label{training-results}

\textbf{Dataset}: 119,384 bookings (89,538 train / 29,846 test) with 50
engineered features and chronological splitting.

\begin{Shaded}
\begin{Highlighting}[]
\CommentTok{\# Load and display training results}
\ImportTok{import}\NormalTok{ json}
\ControlFlowTok{with} \BuiltInTok{open}\NormalTok{(}\StringTok{\textquotesingle{}../artifacts/metrics.json\textquotesingle{}}\NormalTok{) }\ImportTok{as}\NormalTok{ f:}
\NormalTok{    results }\OperatorTok{=}\NormalTok{ json.load(f)}

\CommentTok{\# Create performance summary}
\NormalTok{performance\_data }\OperatorTok{=}\NormalTok{ \{}
    \StringTok{\textquotesingle{}Metric\textquotesingle{}}\NormalTok{: [}\StringTok{\textquotesingle{}Accuracy\textquotesingle{}}\NormalTok{, }\StringTok{\textquotesingle{}ROC{-}AUC\textquotesingle{}}\NormalTok{, }\StringTok{\textquotesingle{}PR{-}AUC\textquotesingle{}}\NormalTok{],}
    \StringTok{\textquotesingle{}Score\textquotesingle{}}\NormalTok{: [results[}\StringTok{\textquotesingle{}accuracy\textquotesingle{}}\NormalTok{], results[}\StringTok{\textquotesingle{}roc\_auc\textquotesingle{}}\NormalTok{], results[}\StringTok{\textquotesingle{}pr\_auc\textquotesingle{}}\NormalTok{]],}
    \StringTok{\textquotesingle{}Interpretation\textquotesingle{}}\NormalTok{: [}
        \StringTok{\textquotesingle{}Overall prediction correctness\textquotesingle{}}\NormalTok{,}
        \StringTok{\textquotesingle{}Ranking quality (all classes)\textquotesingle{}}\NormalTok{, }
        \StringTok{\textquotesingle{}Performance on cancellations (minority class)\textquotesingle{}}
\NormalTok{    ]}
\NormalTok{\}}
\DecValTok{1}
\ImportTok{import}\NormalTok{ pandas }\ImportTok{as}\NormalTok{ pd}
\NormalTok{performance\_df }\OperatorTok{=}\NormalTok{ pd.DataFrame(performance\_data)}
\NormalTok{display(performance\_df.style.}\BuiltInTok{format}\NormalTok{(\{}\StringTok{\textquotesingle{}Score\textquotesingle{}}\NormalTok{: }\StringTok{\textquotesingle{}}\SpecialCharTok{\{:.3f\}}\StringTok{\textquotesingle{}}\NormalTok{\}))}
\end{Highlighting}
\end{Shaded}

\begin{longtable}[]{@{}llll@{}}
\caption{}\label{T_ea93a}\tabularnewline
\toprule\noalign{}
~ & Metric & Score & Interpretation \\
\midrule\noalign{}
\endfirsthead
\toprule\noalign{}
~ & Metric & Score & Interpretation \\
\midrule\noalign{}
\endhead
\bottomrule\noalign{}
\endlastfoot
0 & Accuracy & 0.820 & Overall prediction correctness \\
1 & ROC-AUC & 0.899 & Ranking quality (all classes) \\
2 & PR-AUC & 0.866 & Performance on cancellations (minority class) \\
\end{longtable}

\subsubsection{Final Results Summary}\label{final-results-summary}

\textbf{Best Model: Random Forest (300 estimators)} - \textbf{ROC-AUC:
0.918} (+9.1\% vs logistic regression) - \textbf{PR-AUC: 0.895} (+9.7\%
vs logistic regression)\\
- \textbf{Accuracy: 84.9\%}

\textbf{Key Finding}: Tree models significantly outperform linear
baselines by capturing non-linear cancellation patterns and feature
interactions that logistic regression cannot represent.

\begin{Shaded}
\begin{Highlighting}[]
\CommentTok{\# Model Performance Summary {-} All 11 Experiments}
\NormalTok{df\_results }\OperatorTok{=}\NormalTok{ pd.read\_csv(}\StringTok{\textquotesingle{}../artifacts/results.csv\textquotesingle{}}\NormalTok{)}

\CommentTok{\# Create concise summary table}
\NormalTok{summary\_data }\OperatorTok{=}\NormalTok{ []}
\ControlFlowTok{for}\NormalTok{ \_, row }\KeywordTok{in}\NormalTok{ df\_results.iterrows():}
\NormalTok{    model\_args }\OperatorTok{=} \BuiltInTok{eval}\NormalTok{(row[}\StringTok{\textquotesingle{}model\_args\textquotesingle{}}\NormalTok{]) }\ControlFlowTok{if}\NormalTok{ row[}\StringTok{\textquotesingle{}model\_args\textquotesingle{}}\NormalTok{] }\OperatorTok{!=} \StringTok{\textquotesingle{}}\SpecialCharTok{\{\}}\StringTok{\textquotesingle{}} \ControlFlowTok{else}\NormalTok{ \{\}}
    
    \ControlFlowTok{if}\NormalTok{ row[}\StringTok{\textquotesingle{}model\textquotesingle{}}\NormalTok{] }\KeywordTok{in}\NormalTok{ [}\StringTok{\textquotesingle{}logreg\_baseline\textquotesingle{}}\NormalTok{, }\StringTok{\textquotesingle{}logreg\_l2\textquotesingle{}}\NormalTok{, }\StringTok{\textquotesingle{}logreg\_l1\textquotesingle{}}\NormalTok{]:}
\NormalTok{        model\_family }\OperatorTok{=} \StringTok{\textquotesingle{}Logistic Regression\textquotesingle{}}
    \ControlFlowTok{elif}\NormalTok{ row[}\StringTok{\textquotesingle{}model\textquotesingle{}}\NormalTok{] }\OperatorTok{==} \StringTok{\textquotesingle{}random\_forest\textquotesingle{}}\NormalTok{:}
\NormalTok{        model\_family }\OperatorTok{=} \StringTok{\textquotesingle{}Random Forest\textquotesingle{}}  
    \ControlFlowTok{elif}\NormalTok{ row[}\StringTok{\textquotesingle{}model\textquotesingle{}}\NormalTok{] }\OperatorTok{==} \StringTok{\textquotesingle{}xgboost\textquotesingle{}}\NormalTok{:}
\NormalTok{        model\_family }\OperatorTok{=} \StringTok{\textquotesingle{}XGBoost\textquotesingle{}}
    
\NormalTok{    summary\_data.append(\{}
        \StringTok{\textquotesingle{}Model Family\textquotesingle{}}\NormalTok{: model\_family,}
        \StringTok{\textquotesingle{}ROC{-}AUC\textquotesingle{}}\NormalTok{: row[}\StringTok{\textquotesingle{}roc\_auc\textquotesingle{}}\NormalTok{],}
        \StringTok{\textquotesingle{}PR{-}AUC\textquotesingle{}}\NormalTok{: row[}\StringTok{\textquotesingle{}pr\_auc\textquotesingle{}}\NormalTok{]}
\NormalTok{    \})}

\CommentTok{\# Get best performance per family}
\NormalTok{summary\_df }\OperatorTok{=}\NormalTok{ pd.DataFrame(summary\_data)}
\NormalTok{best\_performance }\OperatorTok{=}\NormalTok{ summary\_df.groupby(}\StringTok{\textquotesingle{}Model Family\textquotesingle{}}\NormalTok{).agg(\{}
    \StringTok{\textquotesingle{}ROC{-}AUC\textquotesingle{}}\NormalTok{: }\StringTok{\textquotesingle{}max\textquotesingle{}}\NormalTok{,}
    \StringTok{\textquotesingle{}PR{-}AUC\textquotesingle{}}\NormalTok{: }\StringTok{\textquotesingle{}max\textquotesingle{}}
\NormalTok{\}).}\BuiltInTok{round}\NormalTok{(}\DecValTok{3}\NormalTok{)}

\BuiltInTok{print}\NormalTok{(}\StringTok{"📊 Best Performance by Model Family:"}\NormalTok{)}
\NormalTok{display(best\_performance)}
\end{Highlighting}
\end{Shaded}

\begin{verbatim}
📊 Best Performance by Model Family:
\end{verbatim}

\begin{longtable}[]{@{}lll@{}}
\toprule\noalign{}
& ROC-AUC & PR-AUC \\
Model Family & & \\
\midrule\noalign{}
\endhead
\bottomrule\noalign{}
\endlastfoot
Logistic Regression & 0.841 & 0.798 \\
Random Forest & 0.918 & 0.895 \\
XGBoost & 0.899 & 0.866 \\
\end{longtable}

\textbf{Model Comparison Results:}

Random Forest achieves the best performance across all metrics, with
XGBoost as a strong second choice and logistic regression providing an
interpretable baseline. The performance gap demonstrates the value of
tree-based methods for capturing complex cancellation patterns.

\subsubsection{Production
Recommendation}\label{production-recommendation}

\textbf{Deploy Random Forest (100 estimators)} for optimal balance of
performance (ROC-AUC: 0.917) and efficiency, with tree-based models
showing +9\% improvement over linear baselines through superior capture
of non-linear cancellation patterns.

\section{Evaluation}\label{evaluation}

\subsection{Model Evaluation}\label{model-evaluation}

\textbf{Performance Validation}: Chronological train/test split on
29,846 test bookings ensures realistic out-of-time validation. Random
Forest achieves ROC-AUC=0.917 and PR-AUC=0.895, representing excellent
discrimination and precision-recall performance for the 37\%
cancellation imbalance.

\textbf{Business Impact}: +9\% performance improvement enables
significantly better booking risk stratification for revenue management
and targeted retention strategies.

\subsubsection{Key Modeling Decisions}\label{key-modeling-decisions}

\textbf{Critical choices}: Chronological split (prevents temporal
leakage), class weighting (handles 37\% imbalance), feature
standardization for logistic regression only, Random Forest selection
(captures non-linear patterns).

\textbf{Limitations}: 2015-2017 Portuguese data may not generalize to
other regions/periods. Monitor for performance drift and feature
importance shifts in production.

\section{Conclusion and Next Steps}\label{conclusion-and-next-steps}

\subsection{Project Summary}\label{project-summary}

Successfully developed a hotel cancellation risk prediction system using
119,384 Portuguese hotel bookings (2015-2017). Random Forest achieved
ROC-AUC=0.917 (+9\% vs logistic regression) through superior capture of
non-linear patterns in lead time, market segments, and customer
commitment signals.

\textbf{Key Business Insights:} - Lead time is the strongest predictor
with clear risk thresholds - Market segment × distribution channel
interactions reveal distinct risk patterns\\
- Customer commitment signals (special requests, parking) indicate
booking stability - Tree models significantly outperform linear
approaches for this domain

\subsection{Recommended Next Steps}\label{recommended-next-steps}

\begin{itemize}
\tightlist
\item
  Deploy Random Forest model in production API for real-time risk
  scoring
\item
  Implement monitoring dashboard with performance drift detection
\item
  Integrate with overbooking and targeted retention systems
\end{itemize}

\textbf{Success Metrics:} - Maintain ROC-AUC ≥ 0.91 and PR-AUC ≥ 0.89 -
Measure revenue impact through RevPAR improvement - API latency ≤ 100ms
with 99.9\% uptime

This system provides actionable booking risk assessment for revenue
optimization, moving from proof-of-concept to production-ready business
tool.

\section{References}\label{references}

The canonical BibTeX entry for this dataset is stored in
\texttt{references.bib} (key:
\texttt{jessemostipak\_hotel\_booking\_demand}).

\phantomsection\label{refs}
\begin{CSLReferences}{1}{0}
\bibitem[\citeproctext]{ref-antonio_almeida_nunes_2019}
Antonio, N., A. de Almeida, and L. Nunes. 2019. {``Hotel Booking Demand
Datasets.''} \emph{Data in Brief} 22: 41--49.
\url{https://doi.org/10.1016/j.dib.2018.11.126}.

\bibitem[\citeproctext]{ref-jessemostipak_hotel_booking_demand}
Mostipak, J. 2020. {``Hotel Booking Demand.''} Kaggle dataset.
\url{https://www.kaggle.com/datasets/jessemostipak/hotel-booking-demand}.

\end{CSLReferences}




\end{document}
